% id: 54-e7
% book: 개념원리 미적분Ⅰ · 06 연속함수의 성질
% section: 필수·발전 예제
% figure: none
% answer: ⑴ 최댓값: $4$, 최솟값: $-5$ ⑵ 최댓값: $2$, 최솟값: $\dfrac{2}{3}$
% answer_source: 본문 풀이
% variant_level: 0 (원본 전사)
% 이 파일은 build.mjs 가 items.json 에서 만든다. 고칠 때는 items.json 을 고친다.
\smash{\raisebox{1.5mm}{\dmkichul{개념원리 미적분Ⅰ 54쪽 예제 7 · 필수}}}
주어진 구간에서 다음 함수 $f(x)$의 최댓값과 최솟값을 구하시오.
\dmsubprob{1} $f(x)=-x^2+2x+3$ \quad $[-2{,}\mkern3mu\mathopen{}2]$
\dmsubprob{2} $f(x)=\dfrac{2}{x-2}$ \quad $[3{,}\mkern3mu\mathopen{}5]$
